% !TeX root = ../教材题库_全部条目_转写工作稿.tex
% Source: XJ-B1-C05-RV-P209-Y1208
% 第5章 / 章末复习 复习题五 / 教材 p.209 / PDF 第216页
% OCR未校对：公式与文字均可能有误；未完成数学TeX转写；禁止直接发给学生
% 原题及相邻原解截图（核对用）：../题目原图/XJ-B1-C05-RV-P209-Y1208-01.webp
% 本文件只在转写完成后作为题干使用；OCR 条目尚不是可靠题干。
\begin{bankitem}{复习题五 第24题}{第5章 / 章末复习 复习题五 / 教材 p.209 / PDF 第216页}{XJ-B1-C05-RV-P209-Y1208}{OCR未校对：公式与文字均可能有误；未完成数学TeX转写；禁止直接发给学生}
\begin{ocrdraft}
[教材p.209 / PDF第216页]
\par
24．如图，A与P分别是单位圆0上的定点与动点，角x
\par
的始边为射线OA，终边为射线OP，过点P作直线OA的垂
\par
线，垂足为M，将点M到直线OP的距离MN表示为的函
\par
数丆0）．
\par
(1)当工e[0，2丌）时，求出函数不到的表达式；
\par
（第24题）
\par
（2）当xeR时，利用计算机软件作出函数3，一丆0）与g(x)一sin2的图象，
\par
通过图象讨论上述两个函数在单调性、最值、奇偶性以及周期性等方面有何
\par
异同．
\end{ocrdraft}
\end{bankitem}
